A cytoskeletal stimulation unit can be conceived as a localized acoustic pad tuned to approximately 100\,Hz, a frequency range associated in the source material with neurite growth and cytoskeletal reorganization. In the RBM framework, the purpose of such a device would be to deliver gentle, repeatable mechanical stimulation to tissue in order to support recovery processes and potentially bias cellular dynamics toward repair-oriented states.

The design premise is straightforward: if cytoskeletal elements participate in mechanically sensitive signaling and structural remodeling, then a focused acoustic interface may help couple external vibration into the tissue microenvironment. A pad form factor is especially suitable for this role because it can maintain stable contact over a defined anatomical region, allowing the stimulus to be applied consistently and with low user burden. In practical terms, the unit would emphasize comfort, controlled intensity, and session repeatability rather than high amplitude output.

Within the broader prototype set, the cytoskeletal stimulation unit complements the vagal patch, sleep sound pad, and SomaTune strap by targeting a more local, tissue-level effect. Its intended use is not generalized relaxation alone, but a recovery-oriented intervention aimed at supporting neurotrophic and structural processes through a 100\,Hz acoustic field.
